\section{問題設定と方法}\label{sec:method}

\subsection{種の価値}

重み $w_{ij}$ を持つグラフ $G=(V,E)$，$|V|=n$ の MAX-CUT を考える．解はスピン $s\in\{-1,+1\}^n$ で表し，カット値を $c(s)=\sum_{(i,j)\in E} w_{ij}[s_i\neq s_j]$ とする．全反転 $s\mapsto -s$ は同じ解とみなす．

CIM の 1 回の実行が出力する解を\textbf{候補}（種）と呼ぶ．候補 $s$ を初期解として後段の探索 $\mathcal{A}$（本論文では主に SA）を予算 $B$ で 1 回実行したときの最良カット値を確率変数 $Y_{\mathcal{A},B}(s)$ とし，その期待値
\begin{equation}
  V_{\mathcal{A},B}(s) = \mathbb{E}\left[Y_{\mathcal{A},B}(s)\right]
\end{equation}
を候補の\textbf{価値}と呼ぶ．期待値は後段の探索の乱数についてとる．本論文の問いは次の 3 つである．

\begin{description}
  \item[C1（現象）] 同じカット値の候補でも $V$ が異なるか．
  \item[C2（説明）] その差を，カット値以外の安価な指標で予測できるか．
  \item[C3（実用）] 指標で候補を選ぶと，CIM の生成を含めた同じ総時間で到達品質が上がるか．
\end{description}

\subsection{CIM と後段の探索}

\paragraph{CIM.} 位相感応増幅器と測定フィードバックを持つファイバーループ型 CIM の進行波モデル\cite{inoue2022}を用いる（CIM 全般は\cite{wang2013,inagaki2016,mcmahon2016}）．ポンプ強度をラウンドごとに一定量（\texttt{dP\_per\_round}）ずつ上げ，各ラウンドで結合 $J=cA$（$A$ は重み付き隣接行列，$c<0$）を通したフィードバックと真空雑音を加える．各実行で，途中で最もカット値の高かった符号を出力とする．パラメータは，G22 について以前の研究で調整した物理定数を用い，問題ごとに結合の強さ $c$ だけを CIM 単体の平均カット値が最大になるように選んだ（後段のラベルは見ていない）．

\paragraph{SA.} 後段の探索は，1 反転提案ごとに温度を指数的に下げる SA\cite{kirkpatrick1983}である．初期解 $s$ から出発し，途中の最良状態（初期解を含む）を返す．予算は反転提案の総数で表し，頂点数で割った値を\textbf{掃引}と呼ぶ．開始温度 $T_0$ は，warm start 向けの $0.5$，冷却 SA の既定値 $2.0$，以前の研究で G22 向けに調整した $4.21$ を用いた．終了温度は $0.001$（調整済みの設定のみ $0.052$）である．

\paragraph{ラベル.} 候補ごとに SA を $R$ 回（互いに異なる乱数）実行し，最良カット値の平均 $\bar{y}(s)$ を $V(s)$ の推定値（ラベル）とする．小規模問題では，SA が最適値に到達した割合をラベルとした．

\subsection{同じカット値の候補の比較}

CIM の出力は 1 反転の局所最適解とは限らず，貪欲な 1 反転局所探索（LS）でただちに改善できることが多い．この「すぐ取れる改善」で説明できる差は自明なので，C1・C2 では，ラベルを初期カット値 $c(s)$ と LS 後のカット値 $c_{\mathrm{LS}}(s)$ で回帰した残差
\begin{equation}
  r(s) = \bar{y}(s) - \left(\beta_0 + \beta_1 c(s) + \beta_2 c_{\mathrm{LS}}(s)\right)
\end{equation}
を対象にする（係数は問題・条件ごとの最小二乗推定）．

\paragraph{C1 の判定.} $R$ 回の SA を偶数回目と奇数回目の 2 群に分け，各群の平均で残差を作り，その相関に Spearman--Brown 補正をかけた値（残差の split-half 信頼性）を求める．SA の乱数は 2 群で独立なので，この値は「カット値と LS で説明できない，候補に固有の価値」がどれだけ再現するかを表す．0.3 以上かつ候補ブートストラップの 95\% 区間の下限が 0 を超えれば，差があると判定した．

\paragraph{C2 の判定.} 指標と残差 $r$ の Spearman 順位相関 $\rho$ と，候補ブートストラップ（各標本で回帰し直す）の 95\% 区間を求める．小規模問題では，候補を（問題，初期カット値）の層に分け，層内の Kendall $\tau_b$ を候補数で重み付け平均した値と，層内で平均を引いた値どうしの相関の 2 乗（偏 $R^2$）を用いた．区間は問題を単位にしたブートストラップで求めた．

\subsection{指標}

\paragraph{固有ベクトル指標.} CIM の増幅は，しきい値の近くでは線形写像 $c_{k+1}\approx g(\sqrt{\eta}I+J)c_k+\xi$ で近似でき，$J$ の最大固有値の固有ベクトルが最初に成長する．$J=cA$（$c<0$）なので，これは $A$ の最小固有値側の固有ベクトルにあたる．$A$ の最小固有値側の 2 本を $u_1,u_2$（固有値の昇順）とし，
\begin{equation}
  a_j(s) = \frac{\left|\langle s, \operatorname{sign}(u_j)\rangle\right|}{n},\qquad E(s) = a_1(s) - a_2(s)
\end{equation}
を定義する．$a_j$ は候補が $u_j$ の符号パターンにどれだけ揃っているかを表す．MAX-CUT と $A$ の最小固有値の関係は Trevisan\cite{trevisan2012}がスペクトル的な丸めとして論じている．固有ベクトルは問題ごとに 1 回だけ計算すればよく，各候補では内積だけで評価できる．

\paragraph{局所的な指標.} 各頂点を反転したときのカットの増分 $g_i$ から，$g_i=0$ の頂点数（\texttt{n\_zero\_gain}），$-2\le g_i\le 0$ の頂点数（\texttt{n\_gain\_minus2\_to\_0}），$g_i=0$ の頂点がつくる連結成分の最大サイズなどを求める．LS 後の解についても同じ量を求める（接頭辞 \texttt{ls\_}）．いずれも辺数に比例する時間で計算できる．

\paragraph{試し打ち.} 候補に短い SA（大規模問題では 50 掃引を 2 回，小規模問題では 1 掃引を 8 回）をかけた結果の平均を指標とする．計算費用は大きいが，候補の後段での振る舞いを直接見る強い比較相手である．

\subsection{厳密量（小規模問題）}

18〜24 頂点の問題では，全反転を同一視した $2^{n-1}$ 状態を全列挙し，次の量を厳密に求めた．
\begin{description}
  \item[障壁] $B(s)=c(s)-\ell^*(s)$．$\ell^*(s)$ は，カット値が $\ell$ 以上の状態だけを 1 反転でたどって $s$ から最適解に行ける最大の $\ell$．$B=0$ ならカット値を下げずに最適解に行ける．状態をカット値の高い順に加えながら素集合森で連結成分を管理して求めた．
  \item[経路長] カット値が $\ell^*(s)$ 以上の状態だけを通る，$s$ から最適解までの最短の反転回数．最適解からの幅優先探索で求めた．
  \item[最適解までの距離] 最も近い最適解とのハミング距離．
\end{description}
全列挙のカット値と，障壁・経路長は，8〜9 頂点の問題で素朴な全探索の結果と一致することを確かめた．
